\subsection{A smaller local input and an exact mixed-prime obstruction}

\begin{statement}{A.46}{Residual conductor control outside a finite explicit set}{Proved}
Let $\ell\ge5$ be prime and let $s$ be a full section of $U_\ell/\Q$,
without a Selmer hypothesis. Put
\[
 T_\ell=\{2,\ell\}\cup\{q:\ q\text{ prime},\ q\mid\ell^2-1\}.
\]
For every $q\notin T_\ell$, the residual Legendre representation
$\bar\rho_{s,\ell}|_{G_{\Q_q}}$ becomes unramified after a quadratic twist.
Consequently, to obtain the reducibility and semisimplification conclusion
of A.6, it suffices that $s$ be point-induced or cuspidal locally at the
finitely many primes in $T_\ell$. No local geometricity assumption at
the other finite primes is needed for this residual-modularity step.
\end{statement}
\begin{proof}
The determinant of the Legendre representation is $\bar\chi_\ell$.
Its geometric image is contained in $\mathrm{SL}_2(\F_\ell)$, of order
$\ell(\ell^2-1)$. Thus, when $q\notin T_\ell$, its geometric restriction
factors through the maximal prime-to-$q$ quotient $P_q$ of the geometric
fundamental group. The resulting arithmetic representation factors
through the corresponding geometrically prime-to-$q$ extension: the
kernel defining that extension is killed by the representation and is
normal in the arithmetic group.

Since $q\ne\ell$, $U_\ell$ has a smooth proper model with disjoint etale
boundary over $\Z_q$. Choose an integral tangent at an $X=0$ cusp.
This gives an unramified genuine action on $P_q$. At this tangent the
Legendre representation is also unramified at $q$: its splitting
quadratic character is unramified at an odd prime, and the Tate Kummer
constant $2^{-8}$ is a $q$-adic unit. We may therefore apply A.28 to the
geometric inertia image of $s$.

If that image is trivial, the residual representation agrees on inertia
with the unramified reference splitting, and is unramified. Otherwise
A.28 puts the prime-to-$q$ section in one cusp decomposition group.
This ramification criterion is also
\citet[Appendix A, Proposition A.6]{AMHoshi2012}, with the prime set
equal to all primes other than $q$; its good-reduction hypothesis is
satisfied here. The calculation A.3 applies at each cusp: after the
local quadratic splitting twist its extension class is
\[
 -8\kappa_\ell(2)+2\ell\kappa_\ell(u)
 =-8\kappa_\ell(2)\quad\text{modulo }\ell.
\]
The Kummer class of a unit is unramified over $\Q_q$ when $q\ne\ell$.
Hence the twisted residual representation is unramified. The splitting
twist itself can ramify, in particular at the infinity cusp, so the
quadratic-twist qualification is retained.

At the primes in $T_\ell$, the assumed point-or-cusp localization gives
exactly A.5's local conclusions: after quadratic twists, unramifiedness
away from $2\ell$, finite flatness and weight two at $\ell$, and conductor
exponent one at two. The residual representation has finite image and
therefore ramifies at finitely many primes. At the remaining primes the
quadratic inertia characters just constructed are thus nontrivial at
only finitely many places. Patch all these characters globally as in
A.6. The resulting odd representation has the same weight-two,
level-two irreducibility contradiction and the same reducible
semisimplification $1\oplus\bar\chi_\ell$. This repeats A.6's argument
with the reduced local input and proves the last assertion.

For example $T_5=\{2,3,5\}$ and $T_7=\{2,3,7\}$. This statement
reduces the local input for R3. The descent requirement in A.30 remains
a separate hypothesis; finite support at density-one places does not
by itself produce an all-place finite-descent adelic point.
\end{proof}



\begin{statement}{A.47}{Complete prime-to-three data can miss wild residual seven inertia}{Proved}
Let $\Pi=\pi_1(U_7/\Q)$, let $P$ be its geometric maximal prime-to-three
quotient, and let $\Pi_P$ be the resulting arithmetic extension. Form the
natural joint quotient of $\Pi$ retaining both $\Pi_P$ and the residual
Legendre representation modulo seven. It has a global section whose
projection to $\Pi_P$ is a fixed global cusp section, but whose residual
Legendre representation has wild inertia of order three at three.
No quadratic twist removes that inertia.

Thus even the entire prime-to-three section, and hence all of its
pro-$\lambda$ projections for $\lambda\ne3$, does not determine this
residual conductor in the natural joint quotient. This is a statement
about that quotient; it does not assert a lift to the full $\Pi$.
\end{statement}
\begin{proof}
The geometric Legendre monodromy modulo seven is
$S=\mathrm{SL}_2(\F_7)$. Indeed its cusp monodromies include
\[
 U=\begin{pmatrix}1&2\\0&1\end{pmatrix},\qquad
 L=\begin{pmatrix}1&0\\-2&1\end{pmatrix},
\]
which generate all upper and lower elementary matrices, and therefore
$S$. These are the standard Legendre monodromy matrices; see
\citet[\S5]{AMStixBirational2014}. Restricting from the tripod to
$U_7$ preserves this image. The geometric subgroup defining the Fermat
cover contains the closed commutator subgroup of the tripod group,
while $S$ is perfect, so the image of that subgroup is still $S$.

We verify both perfectness and the needed prime-to-three assertion
directly. Set $D=\operatorname{diag}(2,4)\in S$, of order three.
For the standard elementary matrices $U(a),L(a)$,
\[
 [D,U(a)]=U(3a),\qquad [D,L(a)]=L(a)
 \quad(a\in\F_7).
\]
The coefficients $3$ and $1$ are nonzero. Thus the commutator subgroup
contains all elementary matrices, proving perfectness. In any quotient
of order prime to three, $D$ is killed; the same identities then kill
every elementary matrix. Consequently $S$ has no nontrivial
prime-to-three quotient. The program
\texttt{a\_main\_local\_matrix.py} independently enumerates the group
of order 336 and verifies these assertions and the normal-closure
calculation.

The geometric image in $P\times S$ is the full product. To see this,
it is a closed subdirect product, so the profinite Goursat argument
identifies a common quotient of $P$ and $S$. Such a quotient is finite
of order prime to three, and hence trivial by the previous paragraph.
The residual arithmetic quotient is
\[
 R=\{(M,g)\in\mathrm{GL}_2(\F_7)\times G_\Q:
                  \det M=\bar\chi_7(g)\}.
\]
Surjectivity onto $G_\Q$ and the full geometric image $S$ prove this
identification. The joint quotient is consequently the entire fibre
product $\Pi_P\times_{G_\Q}R$.

Let $\theta:G_\Q\to\{1,2,4\}\subset\F_7^\times$ be a nontrivial
character of the cyclic cubic extension $\Q(\zeta_9)^+/\Q$.
It has conductor nine: the corresponding order-three character modulo
nine cannot factor through the order-two group modulo three.
Define
\[
 \rho(g)=
 \begin{pmatrix}
 \bar\chi_7(g)\theta(g)&0\\
 0&\theta(g)^{-1}
 \end{pmatrix}.
\]
Its determinant is $\bar\chi_7$, so it defines a section of $R$.
Choose any global cusp section $\sigma$ of $\Pi_P$, for instance the
image of the integral tangent at an $X=0$ cusp. Then
\[
 g\longmapsto\bigl(\sigma(g),(\rho(g),g)\bigr)
\]
is a section of the natural joint quotient just identified.

At three the character $\bar\chi_7$ is unramified, whereas $\theta$
is wildly ramified of order three. Thus a generator of its wild image
acts by $D=\operatorname{diag}(2,4)$. A quadratic character is trivial
on the wild pro-three inertia group, so no quadratic twist removes
this action. In fact each of the two characters has Artin conductor
exponent two at three, giving residual conductor exponent four.
One may choose an arithmetic Frobenius with matrix
$\operatorname{diag}(3,1)$, commuting with $D$ and having determinant
three. All the finite matrix claims are checked by the supplied program.
It also checks the cubic polynomial $x^3-3x+1$ of
$\Q(\zeta_9)^+$: its discriminant is $81$ and the shift
$x^3-3x^2+3$ is Eisenstein at three.

The prime-to-three projection of this section is exactly $\sigma$,
not merely its abelianization. Its wild residual behavior therefore
records a concrete piece of information absent from those projections.
The kernel discarded in forming the joint quotient can obstruct a
further lift. No conclusion about such a lift is used here.
\end{proof}

\begin{statement}{A.48}{Ramification at a bad prime requires its definition}{Proved}
For a geometrically pro-$\Sigma$ section over a local field of residue
characteristic $q\notin\Sigma$, the intrinsic condition
\[
 s(I_q)\text{ acts trivially on the geometric group}
\]
agrees, at good reduction, with trivial geometric inertia in the
unramified model splitting. At bad reduction these descriptions cannot
be interchanged. In particular, if $q\mid n$ with $n$ odd, the intrinsic
inertia action for $U_n/\Q_q$ is already nontrivial on the boundary
part of geometric abelianization, for every nonempty prime set
$\Sigma$ excluding $q$, independently of the section.
\end{statement}
\begin{proof}
The intrinsic definition is
\citet[Appendix A, Definition A.1]{AMHoshi2012}: the composite
$I_q\to G_{\Q_q}\xrightarrow{s}\Pi^\Sigma\to
\operatorname{Aut}(\Delta^\Sigma)$ must be trivial.
His Proposition A.4 identifies this, under good reduction and
$q\notin\Sigma$, with zero geometric inertia after specialization;
Proposition A.6 then recognizes nonzero inertia at a cusp. His
Proposition A.2(ii) says that bad reduction precludes intrinsic
unramifiedness. For an interior point section, Proposition A.7
expresses unramifiedness as good reduction of the curve with that
point also removed. These statements have their indicated hypotheses;
Proposition A.6 does not remove the good-reduction hypothesis from them.

In the Fermat case the relevant nontrivial outer action can be seen
without invoking a model classification. Write $D=X_n\setminus U_n$.
The boundary above $X=0$ contains the $n$ distinct points obtained from
the $n$th roots of unity, up to a fixed rational sign. If $q\mid n$,
inertia at $q$ permutes the $q$th roots nontrivially because
$\Q_q(\mu_q)/\Q_q$ is ramified of degree $q-1>1$.
For any $\lambda\in\Sigma$, the boundary submodule of geometric
abelianization is
\[
 \left(\bigoplus_{d\in D(\overline\Q_q)}\Z_\lambda(1)\right)
       \big/\Z_\lambda(1),
\]
where the denominator is the diagonal relation. The inertia action on
$\Z_\lambda(1)$ is trivial since $\lambda\ne q$, but its permutation
of the boundary basis is not trivial on this quotient: with at least
three boundary points, a nonidentity permutation does not become the
identity modulo the single diagonal relation. It is the kernel of the surjection from geometric pro-$\lambda$
abelianization to $T_\lambda\operatorname{Jac}(X_n)$. Inner conjugation acts
trivially on abelianization. The nontrivial inertia action is therefore
independent of which section supplies a representative of the outer
Galois action.

If instead one fixes a reference section and calls its section cocycle
``unramified'', changing that reference changes the cocycle. At bad
reduction the genuine inertia action is also nontrivial, so the cocycle
on inertia is not a homomorphism into a group with trivial inertia
action. The procyclic geometric-inertia argument of A.28 cannot be
applied under that terminology. This is the precise definitional issue
at the bad primes; it does not assert that any interior point section
has been separated from all pro-$\lambda$ cusp packets. Such a
separation would require an additional argument.
\end{proof}

\begin{statement}{A.49}{Density-one local points do not supply all-place descent}{Proved}
Let $E/\Q$ be $y^2+y=x^3-7$, with origin $O$. There is a family of
distinct full sections of $\pi_1(E)\twoheadrightarrow G_\Q$, indexed
by $\Z_5$, such that the zero member is $s_O$ and every nonzero member
is point-induced over $\Q_q$ for every $q\ne5$ and over $\R$, but is
not point-induced over $\Q_5$. At each $q\ne5$ its point can be chosen
in the finite subgroup $E(\Q_q)[5^\infty]$.

Thus point localization at a cofinite, hence density-one, set of places
does not imply the all-place local condition needed by the
Harari--Stix descent step. This is an actual full-section example on
the proper genus-one curve; it does not assert a lift to its punctured
hyperbolic curve.
\end{statement}
\begin{proof}
By A.10 and A.14, $E(\Q)\simeq\Z/3$ and $\Sha(E/\Q)$ is finite.
Put $T=T_5E$, $S=\{3,5,\infty\}$, and
$G_S=\operatorname{Gal}(\Q_S/\Q)$. The conductor is $27$, so the
action on $T$ factors through $G_S$.

We first prove that $H^1(G_S,T)$ has positive $\Z_5$-rank. For
$M_n=E[5^n]$, the global Euler characteristic formula of
\citet[I, Theorem 5.1, p.67]{AMMilneADT2006} gives
\[
 \frac{|H^0(G_S,M_n)|\,|H^2(G_S,M_n)|}
      {|H^1(G_S,M_n)|}
 =\frac{|H^0(\R,M_n)|}{|M_n|}=5^{-n}.
\]
All its hypotheses hold: $S$ is finite and contains the infinite place
and five, and $M_n$ is unramified outside $S$.
Complex conjugation has eigenvalues $1,-1$ on $T$, because its
determinant is cyclotomic and five is odd. Hence the numerator on the
right has order $5^n$, while $|M_n|=5^{2n}$.
Also $H^0(G_S,M_n)=E(\Q)[5^n]=0$ as a group, of order one.

Here is the integral passage, including its finiteness point.
Finite-coefficient cohomology is finite, and its groups above degree
two vanish for these odd-primary coefficients
\citep[I, Theorem 4.10(c), Corollary 4.15]{AMMilneADT2006}.
Inverse limits have no derived-limit term because the relevant systems
are finite and satisfy the Mittag--Leffler condition. The groups
$H^i(G_S,T)$ are compact $\Z_5$-modules. The exact sequences for
$0\to T\xrightarrow{5^n}T\to M_n\to0$ imply that their quotients
modulo five are finite; compact Nakayama therefore makes them finitely
generated. They also give $H^3(G_S,T)=0$ and
$H^1(G_S,T)[5^n]=0$, the latter because both $H^0(G_S,T)$ and
$H^0(G_S,M_n)$ vanish. Write $r_i$ for the rank of $H^i(G_S,T)$.
In
\[
 0\longrightarrow H^i(G_S,T)/5^n
 \longrightarrow H^i(G_S,M_n)
 \longrightarrow H^{i+1}(G_S,T)[5^n]
 \longrightarrow0,
\]
the finite torsion terms in degrees one and two cancel in the Euler
characteristic. The displayed formula consequently yields
$5^{n(r_2-r_1)}=5^{-n}$, so $r_1=r_2+1\ge1$.
Choose a nonzero infinite-order class $z\in H^1(G_S,T)$ and inflate
it to $H^1(\Q,T)$; inflation is injective.

For every finite $q\ne5$, $E(\Q_q)$ is compact with an open pro-$q$
subgroup. Thus its maximal pro-five quotient is the finite subgroup
$E(\Q_q)[5^\infty]$. Local Tate duality
\citep[I, Corollary 3.4, p.43]{AMMilneADT2006} makes
$H^1(\Q_q,E)[5^\infty]$ finite as well. Taking inverse limits in
the local Kummer sequences therefore gives
\[
 H^1(\Q_q,T)=E(\Q_q)^{\wedge5}
            =E(\Q_q)[5^\infty],\qquad q\ne5;
\]
the possible quotient $T_5H^1(\Q_q,E)$ is zero. This argument applies
also at the bad prime three. At the real place
$H^1(\R,T)=0$ because multiplication by two is invertible on $T$.

The all-place compact five-adic Selmer group fits into the inverse
Kummer exact sequence
\[
 0\longrightarrow E(\Q)^{\wedge5}
 \longrightarrow \operatorname{Sel}(T)
 \longrightarrow T_5\Sha(E/\Q)\longrightarrow0.
\]
Both outside terms vanish, so $\operatorname{Sel}(T)=0$.
For every $a\in\Z_5\setminus\{0\}$, the class $az$ is nonzero and
already satisfies the Kummer condition at every place other than five.
Its localization at five is therefore not a Kummer class.

Finally the origin splits the full fundamental extension as
\[
 \pi_1(E)=\left(\prod_\lambda T_\lambda E\right)\rtimes G_\Q.
\]
Insert a continuous cocycle representing $az$ in the five-component
and put zero in every other component. This defines a genuine full
section $s_a$. Distinct $a$ give distinct section classes because their
five-adic cohomology classes are distinct. At $q\ne5$, choose the
point $P_q\in E(\Q_q)[5^\infty]$ giving the five-component. Its
Kummer class in $T_\lambda E$ is zero for every $\lambda\ne5$:
multiplication by every power of $\lambda$ is invertible on the finite
five-primary group. Thus this single point realizes the \emph{entire}
full local section, simultaneously in all coefficient components.
At the real place the section is that of $O$; at five it cannot be
that of a point, since its five-component is not Kummer.
This proves every assertion.
\end{proof}
